\section{Grundlagen}
\label{sec:grundlagen}

\subsection{Multi-Tenant SaaS-Architektur}

In einer Multi-Tenant-Architektur teilen mehrere Mandanten (Organisationen)
dieselbe Applikationsinstanz und, im Shared-Database-Modell, auch dieselbe
Datenbankinstanz. Die logische Trennung der Mandantendaten erfolgt über ein
Diskriminator-Attribut (hier: \texttt{organization\_id}), das in jeder
geschäftsrelevanten Tabelle vorhanden ist. Die Applikationslogik muss bei jedem
Dateienbankzugriff sicherstellen, dass ausschliesslich Datensätze des aktuell
authentifizierten Mandanten ausgegeben werden. Fehlt diese Filterung an einer
zentralen Stelle (etwa in einer gemeinsamen Basisklasse), sind potenziell
alle darauf aufbauenden Endpunkte betroffen.

\subsection{Insecure Direct Object Reference (IDOR)}

IDOR ist laut OWASP \cite{owasp_idor} eine Unterklasse von \emph{Broken Access
      Control} (OWASP Top 10, A01) und in CWE-639 \cite{cwe_639} als
\enquote{Authorization Bypass Through User-Controlled Key} beschrieben.
Bei einem IDOR-Angriff manipuliert ein Angreifer einen direkt
benutzergesteuerten Bezeichner (in diesem Fall eine UUID) um auf Ressourcen
zuzugreifen, ohne eine Autorisierungsprüfung zu durchlaufen.

\textbf{Angriffsszenario auf dnchub-app}
\begin{enumerate}
      \item Angreifer authentifiziert sich als Benutzer von Organisation A.
      \item Über einen normalen API-Aufruf erhält er eine gültige UUID eines
            eigenen Ressourceneintrags (z.\,B. eines Fahrzeugs).
      \item Er ersetzt in einem weiteren Request die UUID durch eine fremde UUID
            (erraten, durch Enumeration oder durch Insiderwissen erhalten).
      \item \texttt{BaseService.get\_or\_404(db, uuid)} fragt die Datenbank mit
            \texttt{WHERE id = uuid} ab, ohne \texttt{organization\_id}-Prüfung.
      \item Die Response enthält die Ressource von Organisation B.
\end{enumerate}

Da alle PATCH- und DELETE-Endpunkte ebenfalls über \texttt{get\_or\_404()}
gehen, kann der Angreifer fremde Datensätze auch verändern oder löschen.

\subsection{Soft-Delete-Muster}

Das Soft-Delete-Muster markiert zu löschende Datensätze mit einem Zeitstempel
(\texttt{deleted\_at}) statt sie physisch zu entfernen. Dies ermöglicht
Audit-Trails und einfache Wiederherstellung. Die Korrektheit des Musters setzt
voraus, dass alle Abfragen konsistent nach \texttt{deleted\_at IS NULL} filtern.
Wird dies, wie vorliegend in \texttt{get()} und \texttt{get\_or\_404()},
unterlassen, sind soft-gelöschte Datensätze weiterhin per direkter ID-Abfrage
zugänglich (GitHub Issue \#14 \cite{github_issue_14}).

\subsection{CVSS und CWE-Klassifikation}

Der Common Vulnerability Scoring System (CVSS) v3.1 \cite{first_cvss31}
bewerte Schwachstellen auf einer Skala von 0 bis 10.
Die Schwachstelle in Issue \#5 erhält einen Score von \textbf{9.3 (Kritisch)}
mit folgender Vektorbewertung:

\begin{itemize}
      \item \textbf{AV:N}, Angriff über das Netzwerk möglich
      \item \textbf{AC:L}, Angriffskomplexität niedrig (keine Sonderbedingungen)
      \item \textbf{PR:L}, niedrige Privilegien erforderlich (nur Login nötig)
      \item \textbf{UI:N}, keine Benutzerinteraktion erforderlich
      \item \textbf{C:H / I:H / A:H}, Hoch: Vertraulichkeit, Integrität und
            Verfügbarkeit aller mandantenübergreifenden Daten betroffen
\end{itemize}
